% A conservative Fourier-mode statement, with grid assumptions explicit.
\clearpage
\section*{Stability laboratory: a smooth picture is not a proof}
\noindent The continuous heat equation on a unit square is
\(\partial_tu=\kappa(\partial_{xx}u+\partial_{yy}u)\).
Its analytical mode
\(u(x,y,t)=e^{-2\kappa\pi^2 t}\sin(\pi x)\sin(\pi y)\)
is smooth and decays. The corresponding finite-difference solver
does not automatically inherit those properties for every time step.

On a uniform grid with spacing \(h\) in both directions, the
forward-time centered-space discretization is
\[
u_{i,j}^{n+1}=u_{i,j}^n+r
(u_{i+1,j}^n+u_{i-1,j}^n+u_{i,j+1}^n+u_{i,j-1}^n-4u_{i,j}^n),
\quad
r=\frac{\kappa\Delta t}{h^2}.
\]
The coefficient \(r\) is dimensionless: with \(\kappa\)
in length-squared per second, \(h^2\) is length-squared and
\(\Delta t\) is seconds. Keeping track of these units catches
many subtle mistakes before numerical analysis begins.

For Fourier modes of the periodic or infinite-grid operator,
the amplification factor is
\[
g(k_x,k_y)=1-4r\left(
\sin^2\frac{k_xh}{2}+\sin^2\frac{k_yh}{2}
\right).
\]
The most demanding mode yields \(g_{\min}=1-8r\).
Requiring \(\lvert g\rvert\le1\) for every Fourier mode
gives \(0\le r\le1/4\), or
\[
\Delta t\le\frac{h^2}{4\kappa}.
\]
This is a classical sufficient limit for the two-dimensional
forward-Euler, centered-difference heat method; boundary details
and the actual finite-grid spectrum must be checked before
transferring a statement about a limiting grid verbatim.

At \(h=1/4\) and \(\kappa=1\), the bound is
\(\Delta t\le0.015625\). Refining \(h\) by a factor of two
requires a fourfold smaller step to preserve the same \(r\).
A more detailed spatial grid can therefore make an explicit
solver much more expensive even when it has fewer visual artifacts.

\paragraph{Reconstruction challenge.}
Implement both forward Euler and backward Euler using the same
discrete Laplacian and boundary constraints.
Record mode amplification, reference solution errors, total energy,
and maximum-principle violations separately.
Show solutions for stable and unstable choices as animations, but
never use a plausible-looking low-frequency image as a stability test:
an unobserved high-frequency mode can still grow.


\paragraph{C++23 checkpoint: do not let a heat plot hide instability.}
The native command \texttt{book02\_workbench heat} exports one row for
each position and time step. The optional
\texttt{tools/labs/book02\_cpp23/plot.py} script only draws those
exported rows; it does not calculate a heat update. With uniform grid
spacing \(\Delta x\), the explicitly stepped one-dimensional heat
equation has the form
\[
u_i^{n+1}=(1-2r)u_i^n+r u_{i-1}^n+r u_{i+1}^n,\qquad
r=\frac{\alpha\,\Delta t}{(\Delta x)^2}.
\]
When \(0\leq r\leq\tfrac12\), each interior update is a convex
combination of its three neighbors. Thus nonnegative input values
stay nonnegative and, with fixed zero endpoints, the discrete maximum
does not grow. The restriction has a physical as well as a computational
interpretation: refining the grid by a factor of two requires reducing
the explicit time step by roughly a factor of four at fixed diffusivity.

The executable rejects \(r>\tfrac12\) before calculating a purported
solution, and keeps its Dirichlet endpoints at zero. To test the
reasoning, start from a symmetric single-cell pulse and derive the
values after one step with \(r=\tfrac14\). They are \(0,\tfrac14,
\tfrac12,\tfrac14,0\). Check symmetry, positivity, and boundary
values independently of the plotting code. Then request
\(r=0.51\): rejection, rather than a pretty but physically misleading
plot, is the expected result.

This is a one-dimensional explicit stencil, not a finite-element
heat-method geodesic solver and not an implicit unconditionally stable
method. The later geometric operators in this book require their own
mass matrices, boundary treatment, linear solves, and convergence study.
